\section{An algebraic central fibre and a transverse coefficient slice}
\label{sec:algebraic-family}

We select algebraic coefficients within the compact eight-node
construction, rather than assume that an unspecified general member
has computable coefficients. We verify every central condition used
by the subsequent period and analytic constructions.

\subsection{The fixed contraction}

Let $\Sigma=\PP^1_f\times\PP^1_g$, $L_0=\OO_\Sigma(1,-1)$,
and use the quotient convention for
\[
 Z=\PP_\Sigma(\OO\oplus L_0)
       \times_\Sigma\PP_\Sigma(\OO\oplus L_0).
\]
Write $\xi_1,\xi_2$ for the tautological classes and $k_i=g+\xi_i$.
These classes are globally generated because their direct images are
$\OO(0,1)\oplus\OO(1,0)$. For base coordinates $s,t$ and local
quotient coordinates $(u_i,v_i)$, the map to $A$ is
\begin{equation}\label{eq:fixed-contraction}
 (x;y)=(v_1s;u_1t),\qquad (x';y')=(v_2s;u_2t).
\end{equation}
Its scheme image is the normal complete-intersection fourfold
\begin{equation}\label{eq:fixed-fourfold}
 V_0=\{f_1=x_0x'_1-x_1x'_0=0,\quad
                 f_2=y_0y'_1-y_1y'_0=0\}.
\end{equation}
Indeed the determinant equations recover the common directions
$s,t$ whenever their corresponding vector pairs are nonzero. A
direction is free exactly when both corresponding pairs vanish;
both cannot be free. This proves irreducibility and the image
description. The equations are a generically smooth regular sequence,
so their complete intersection is reduced, with this scheme image.

The singular locus is the disjoint union
\[
 \ell_x=\{x=x'=0,\ [y]=[y']\},\qquad
 \ell_y=\{y=y'=0,\ [x]=[x']\}.
\]
Near either line one equation eliminates a smooth variable and the
other is a $2\times2$ determinant in four variables. Thus its
transverse singularity is an ordinary threefold node. The complete
intersection is Cohen--Macaulay and regular in codimension one,
which proves normality.

The map $q_Z:Z\to V_0$ contracts a ruling on each of the two
sections where both quotients are onto $\OO$ or both onto $L_0$.
These exceptional sections have codimension two. Adjunction gives
\begin{equation}\label{eq:central-anticanonical}
 -K_Z=4g+2\xi_1+2\xi_2=2k_1+2k_2
             =q_Z^*\OO_{V_0}(2,2),\qquad
 K_{V_0}=\OO_{V_0}(-2,-2).
\end{equation}
The contraction is small and crepant, and normality gives
$q_{Z*}\OO_Z=\OO_{V_0}$. We use the exact polarization
$H=\OO_{V_0}(2,2)|_{X_0}$, also writing $H$ for its first Chern class.

\subsection{Finite admissibility conditions}

Let $\mathcal H=H^0(A,\OO(2,2))\cong\CC^{100}$, and set
$X_h=V_0\cap\{h=0\}$ and $Y_h=q_Z^{-1}X_h$.
The map $\mathcal H\to H^0(Z,-K_Z)$ is surjective: the twisted
Koszul resolution
\[
 0\longrightarrow\OO_A\longrightarrow\OO_A(1,1)^{\oplus2}
            \longrightarrow I_{V_0}(2,2)\longrightarrow0
\]
has no higher cohomology in its first two terms. Normality and
projection formula then give the assertion. The target has dimension
$100-2\cdot16+1=69$.

The fixed toric fan has rays
\begin{equation}\label{eq:toric-rays}
 e_1,-e_1,e_2,-e_2,f_0,e_1+e_2-f_0,g_0,-e_1-e_2-g_0.
\end{equation}
Its sixteen maximal cones choose one ray from each consecutive
pair and are unimodular. Integer linear algebra supplies the
sixteen affine charts and anticanonical trivializations. The full
torus support can be computed directly by substituting
\begin{equation}\label{eq:torus-substitution}
 (x_0,x_1,y_0,y_1)=(as,a,t,1),\qquad
 (x'_0,x'_1,y'_0,y'_1)=(bs,b,t,1)
\end{equation}
in the standard $(2,2)$ monomials. It has 69 distinct monomials.
This avoids a choice of sign for the polar-polytope convention.

Define $\mathcal H^{\rm adm}$ by the following conditions.
\begin{enumerate}
\item In each toric chart, the polynomial $H_\alpha$ representing
      $q_Z^*h$ has no common zero with its four partial derivatives.
\item On each of $\ell_x,\ell_y$, the restricted binary quartic has
      nonzero discriminant and nonzero values at the two fixed endpoints.
\item All coefficients in the 69-monomial support are nonzero. For
      every face $\Theta$, including the full polytope, the face
      polynomial $H_\Theta$ has no common zero in $(\CC^*)^4$ with
      $s\partial_sH_\Theta,t\partial_tH_\Theta,
       a\partial_aH_\Theta,b\partial_bH_\Theta$.
\end{enumerate}
These are nonempty Zariski-open conditions over $\QQ$. The first
is openness of smoothness in a projective divisor family; the
second is a finite set of polynomial nonvanishing conditions.
The third is the face-nondegeneracy condition for the fixed Newton
polytope. The general nondegenerate anticanonical construction in
\cite[Appendix~A]{W_Toric} satisfies smoothness and the quartic
conditions. Surjectivity of the section map permits the additional
nonzero coefficient conditions. Each open condition is dense in the
irreducible space $\mathcal H$, so their finite intersection is nonempty.

\begin{proposition}\label{prop:algebraic-choice}
Membership in $\mathcal H^{\rm adm}$ is decidable for represented
algebraic coefficients. Enumerating $\QQ(\sqrt{-1})^{100}$ and
testing membership therefore terminates with a coefficient vector
$h_*\in\mathcal H^{\rm adm}$.
\end{proposition}

\begin{proof}
For the first condition, test whether each ideal generated by
$H_\alpha$ and its derivatives contains one. For the third, clear
Laurent monomials and adjoin $zstab-1$, with a new variable $z$;
the ideal contains one exactly when the torus zero set is empty.
Gr\"obner reduction over the represented number field and the
Nullstellensatz decide these tests. Quartic discriminants, endpoint
values, support coefficients and polytope faces require finite exact
algebra as well. Every candidate is decided. The rational complex
coefficient vectors are Zariski dense, so one lies in the open set.
\end{proof}

\subsection{Central invariants and attachment}

For every admissible $h$, the divisor $Y_h$ is smooth. Away from
the two singular lines $q_Z$ is an isomorphism, so $X_h$ has
no singularities there. At a simple quartic zero, eliminate the
coordinate along its line; the transverse determinant becomes
exactly a threefold node. Thus $X_h$ has eight nodes, in two
groups of four, and $Y_h\to X_h$ is a projective small crepant
resolution. Its exceptional curves have normal bundle
$\OO(-1)\oplus\OO(-1)$: on the all-$\OO$ section the normal
bundle in $Z$ is $(L_0^{-1})^{\oplus2}$, and on the all-$L_0$
section it is $L_0^{\oplus2}$; the simple quartic derivative
removes the trivial direction normal to the ruling. In divisor
coordinates $(f,g,\xi_1,\xi_2)$, their degree vectors are
\begin{equation}\label{eq:exceptional-types}
          \Gamma_f=(1,0,0,0),\qquad \Gamma_g=(0,1,-1,-1).
\end{equation}

We state the fixed global toric data used here. A nondegenerate
anticanonical member in this smooth toric refinement is simply
connected and has $(h^{1,1},h^{2,1})=(4,60)$; its four indicated
divisor classes span rational $H^2$. The Hodge-number calculation
is \cite[Appendix~A]{W_Toric}, and the lattice-quotient verification
of simple connectivity is \cite[Section~8]{W_HodgeRay}.
These assertions apply to every admissible member: nondegeneracy
is included above, and the smooth proper family over the connected
open set $\mathcal H^{\rm adm}$ also transports the topology from
a general member. No additional Picard-rank genericity is assumed.

For clarity, the relevant intersection calculation is
\begin{equation}\label{eq:central-intersection}
 \left(\int_{Y_h}D_iD_jH\right)_{i,j}
 =\begin{pmatrix}
 0&8&4&4\\8&0&12&12\\4&12&-8&0\\4&12&0&-8
 \end{pmatrix},\qquad (D_i)=(f,g,\xi_1,\xi_2).
\end{equation}
It follows by integrating $D_iD_j(-K_Z)^2$ using
$f^2=g^2=0$, $\xi_i^2=(f-g)\xi_i$, and
$\int_Zfg\xi_1\xi_2=1$. Its determinant is $-16384$ and
$H^3=320$. Thus the four classes are independent, and
\eqref{eq:exceptional-types} has rank two.

Delete disjoint node balls to obtain the common exterior $U$.
Thom and Mayer--Vietoris give, on a nearby smooth fibre,
\begin{equation}\label{eq:integral-attachment}
 H^2(Y_h;\ZZ)\xrightarrow{r}\ZZ^8
      \xrightarrow{\partial}H^3(X;\ZZ),\qquad
 r(D)_i=D\cdot C_i,\quad
 \partial e_i=\operatorname{PD}_L(\delta_i).
\end{equation}
Indeed the exceptional curves have real codimension four, so
$H^2(Y_h)\to H^2(U)$ is an isomorphism. A resolution cap retracts
to $S^2$ and its link restriction records the curve degree; a
smoothing cap retracts to $S^3$ and has no degree-two cohomology.
The connecting image of its link generator is the oriented Thom
dual of that sphere. These local descriptions identify the maps
and give exactness, with the left duality convention used below.

The restrictions of $f,g$ are the two block indicators
$v_f,v_g\in\ZZ^8$. The rational image of $r$ is their span by
\eqref{eq:exceptional-types} and the spanning divisor classes.
Every integral vector in this rational span is constant with an
integer value on each block; hence $\im r=\ZZ v_f\oplus\ZZ v_g$.
It follows that the eight local classes have precisely these
integral relations and rank-six span. The following integral argument
also determines the saturation of their image.

\begin{proposition}\label{prop:primitive-vanishing-lattice}
For every admissible $h$, one has $H_2(Y_h;\ZZ)\cong\ZZ^4$
and $H^3(Y_h;\ZZ)$ is torsion-free. The vanishing lattice
$L=\langle\delta_1,\ldots,\delta_8\rangle\subset H_3(X;\ZZ)$
is primitive, and
\[
                    H_3(X;\ZZ)/L\cong\ZZ^{128}.
\]
\end{proposition}

\begin{proof}
The very ample bundle $\OO(2,2)$ embeds $V_0$ so that $X_h$
is a hyperplane section. Its complement is affine, and
\[
 q_Z:Z\setminus Y_h\longrightarrow V_0\setminus X_h
\]
is proper, being the base change of the projective contraction.
Its fibres are points off the two singular lines and projective
lines over their complements of the four quartic zeros.
If $\phi(k)$ is the complex dimension of the locus of fibres of
dimension $k$, then $\phi(0)=4$ and $\phi(1)=1$, with no other
nonempty fibre loci. The homotopy-dimension theorem of
Goresky--MacPherson \cite[Part~II, \S1.1*, p.~152]{GoreskyMacPherson},
applied to the proper map into the affine hyperplane complement,
gives a CW complex of real dimension at most
$\max_k(2k+\phi(k))=4$ homotopy equivalent to $Z\setminus Y_h$.
Integral duality on the compact real eight-manifold $Z$ gives
\[
 H_k(Z,Y_h;\ZZ)\cong H^{8-k}(Z\setminus Y_h;\ZZ)=0
                         \qquad(k\le3).
\]
The pair sequence therefore identifies $H_2(Y_h;\ZZ)$ with
$H_2(Z;\ZZ)=\ZZ^4$, the latter by the projective-bundle formula.
Universal coefficients now imply that $H^3(Y_h;\ZZ)$ is torsion-free.

Write $N_i\cong S^2\times S^3$ for the node links, $B_i\simeq S^3$
for the smoothing caps, and $B'_i\simeq S^2$ for the resolution caps.
On the smoothing side, Mayer--Vietoris continues
\eqref{eq:integral-attachment} as
\[
 \ZZ^8\xrightarrow{\partial}H^3(X;\ZZ)
 \longrightarrow H^3(U;\ZZ)\oplus\bigoplus_iH^3(B_i;\ZZ)
 \longrightarrow\bigoplus_iH^3(N_i;\ZZ).
\]
Each cap-to-link restriction in degree three is an isomorphism.
The kernel of the last arrow consequently projects isomorphically
onto $H^3(U;\ZZ)$, so $\operatorname{coker}\partial\cong H^3(U;\ZZ)$.
On the resolution side, the degree-two restrictions
$H^2(B'_i;\ZZ)\to H^2(N_i;\ZZ)$ are isomorphisms. Hence the
corresponding sequence gives
\[
 0\longrightarrow H^3(Y_h;\ZZ)\longrightarrow H^3(U;\ZZ)
                    \longrightarrow\bigoplus_iH^3(N_i;\ZZ)=\ZZ^8.
\]
Thus $H^3(U;\ZZ)$ is an extension of torsion-free groups and is
torsion-free. Poincar\'e duality identifies it with
$H_3(X;\ZZ)/L$, whose rank is $134-6=128$.
\end{proof}

Orient the spheres accordingly and put
\begin{equation}\label{eq:period-space}
 R=\left\{p\in\CC^8:
       \sum_{i\in I_f}p_i=\sum_{i\in I_g}p_i=0\right\}.
\end{equation}
The toric anticanonical sequence gives $K_{Y_h}=\OO$ and
$H^1(\OO_{Y_h})=H^2(\OO_{Y_h})=0$. Rationality of the nodes
descends these vanishings to $X_h$, and adjunction gives
$\omega_{X_h}=\OO$.

We use the following precise deformation consequence from
\cite[Section~2]{W_Hodge} together with the relative embedding in
\cite[Section~2]{W_HYM}. For this fixed small transition there is
a proper polarized smoothing over a full ball in $R$, with smooth
total space and a nowhere-zero relative three-form $\Omega$, whose
oriented local periods are its coordinates. Its local equations are
$\sum_{j=1}^4w_j^2=p_i/u_i(p)$, with holomorphic units $u_i$;
its full discriminant is $\bigcup_i\{p_i=0\}$. The family embeds
relatively in $A$ with the three fixed bidegrees and exact polarization
$\OO(2,2)$. These results require the Calabi--Yau vanishings,
projective small resolution, exceptional rank-two map, and full-support
smoothing relation just verified. They follow from unobstructed
polarized deformation, the local-to-global smoothing map, the attachment
sequence, and extension of the two ambient line bundles. The conclusion
is an existence theorem; we do not assume that its arbitrary holomorphic
embedding data are evaluated.

\subsection{Rank and algebraic first derivatives of the period map}

Fix the selected $h_*$ and $c_*=(f_1,f_2,h_*)\in\mathcal C$.
On a sufficiently small coefficient neighbourhood, the complete
intersections are flat and Gorenstein. Fix the standard algebraic
trivialization of
\[
 K_A\otimes\det(\OO(1,1)^{\oplus2}\oplus\OO(2,2))=\OO_A
\]
and the order of the equations. Adjunction gives a single relative
residue generator $\Psi$, including at nodal fibres. Local Morse
coordinates and the sphere integral give holomorphic periods
\begin{equation}\label{eq:coefficient-period-units}
 P_i(c)=\tau_i(c)u_i(c),\qquad u_i(c_*)\ne0.
\end{equation}
The attachment relations hold on the common exterior and its nearby
fillings. They imply $P(c)\in R$ on smooth fibres, and then
everywhere by holomorphic continuation.

Apply the preceding existence theorem to this selected central
member. Its relative embedding supplies a holomorphic coefficient
map $j$ from a ball in $R$ with $j(0)=c_*$. The quotient of
$j^*\Psi$ by $\Omega$ is a holomorphic unit on the total space.
Properness, connected fibres and $H^0(\OO_{X_p})=\CC$ imply
$\pi_*\OO=\OO$, so this quotient is a unit $a(p)$ on the base.
Consequently
\begin{equation}\label{eq:rank-period-map}
             P(j(p))=a(p)p,\qquad DP_{c_*}\circ Dj_0=a(0)I_R.
\end{equation}
Thus $DP_{c_*}$ has rank six. Only existence of $j$ was used;
no numerical inversion of that particular map is required.

The node coordinates are algebraic quartic roots and can be isolated.
At each node choose a nonzero two-equation Jacobian minor and
eliminate its two variables. The first two jets of the remaining
four-variable equation are algebraic. Its central Hessian is
nonsingular; symmetric elimination and algebraic square roots put
its quadratic part into $\sum w_j^2$. The critical value derivative
$\dd\tau_i(c_*)$ is therefore algebraic. If $a_i\ne0$ is the
algebraic central residue coefficient in these coordinates, scaling
the three-sphere gives
\begin{equation}\label{eq:period-first-derivative}
       \dd P_i(c_*)=\epsilon_i\pi^2a_i\,\dd\tau_i(c_*),
                   \qquad\epsilon_i\in\{1,-1\}.
\end{equation}
The unit's leading term uses only the linear Jacobian and residue
value; higher Morse terms do not contribute because the period
contains the factor $\tau_i$. The full bounded version of this
calculation is proved in Section~\ref{sec:local-periods}.

Test the finitely many signs in \eqref{eq:period-first-derivative}
for which the two block row sums vanish, after dividing by $\pi^2$.
These are algebraic equality tests. The attachment relations and
rank six show that a compatible choice exists and that its two
block rows span the entire left kernel. Other successful choices
differ by a common sign on each block. This uses only the known
source relations, not coordinates of any local sphere in the global
middle lattice.

Retain the first three rows in each block to identify $R$ with
$\CC^6$. The resulting $6\times132$ derivative has a nonzero
six-column minor. Removing its common power of $\pi^2$ makes
selection of such a minor an exact algebraic operation. The six
coordinate directions define an affine six-plane through $c_*$,
with algebraic coefficients at algebraic parameter values. This is
the slice required in Theorem~\ref{thm:evaluated-periods}. That
theorem supplies an evaluated inverse on a known full period ball
and a smooth algebraic reference fibre. No ratio restriction on
its nonzero periods has been introduced.
